
Research pipelines can complete successfully while producing scientifically invalid conclusions. All JSON schemas validate, all type checks pass, all modules execute without errors---yet the core scientific claim is contradicted by empirical evidence. Existing pipeline validation tools (MLflow, DVC, Great Expectations) excel at catching syntactic errors but remain blind to reasoning failures because they treat execution traces as opaque logs.

This work addresses the gap by treating Model Context Protocol tool-calling traces as semantic artifacts that encode researcher reasoning in natural language. When a researcher queries a knowledge base with ''search for pruning techniques reducing effective rank,'' the query text encodes an assumption that rank should decrease. When experiments return ''effective rank increased 6.02\%,'' the result text contains evidence. By analyzing both the syntactic structure (tool names, parameter types, result schemas) and semantic content (assumptions in query text, evidence in result text), automated validation becomes possible without manual test specification.

\textbf{The Core Insight.} MCP tool calls embed natural language in both query parameters and result content. Research pipelines are inherently NL-rich because queries are hypothesis-driven and results are findings. This dual-layer encoding---assumptions in what researchers ask, evidence in what they find---enables validation through semantic extraction and comparison.

\textbf{Three Contributions.} First, this work demonstrates that 97.48\% of MCP tool calls contain complete records with natural language in both query parameters and result content (exceeding a 95\% threshold by 2.48 percentage points). This establishes that MCP traces contain sufficient semantic information for dual-layer analysis. Second, LLM-based extraction of assumptions and claims achieves 82.7\% recall and 86.3\% precision with substantial inter-rater agreement (Cohen's $\kappa$=0.716) when validated against independent human annotation, without requiring labeled training data. Third, semantic similarity matching for constraint inference failed to detect assumption-evidence contradictions (0\% recall), establishing a boundary condition: sentence embeddings optimize for topic similarity, not logical contradiction.

\textbf{Refined Claim.} The two-layer framework (syntactic validation and semantic extraction) is empirically validated with 97.48\% trace completeness and 82.7\% recall / 86.3\% precision for extraction. Constraint inference via semantic similarity requires methodological redesign. The full three-layer framework's predicted failure detection rate remains unverified pending integration of entailment models or LLM-based contradiction detection.

\textbf{Organization.} Section 2 reviews related work on MCP frameworks, agent-driven automation, and pipeline validation. Section 3 describes the three-layer architecture and design rationale. Section 4 presents the experimental protocol, dataset, and evaluation metrics. Section 5 reports results for trace completeness, natural language presence, extraction quality, and constraint inference. Section 6 discusses interpretation, limitations, and connections to NLI literature. Section 7 concludes with future directions.
